\documentclass[11pt]{article}
\usepackage[a4paper,margin=25mm]{geometry}
\usepackage{amsmath,amssymb,amsthm}
\usepackage[hidelinks]{hyperref}
\emergencystretch=2em
\newtheorem{theorem}{Theorem}
\newtheorem{lemma}{Lemma}
\title{A balanced P\'olya urn without a positive\\quadratic large-deviation rate\\
\large Disproof of Conjecture 00000002114}
\author{gaochengzhecpu}\date{}
\begin{document}\maketitle
\begin{abstract}
The classical two-colour P\'olya urn with one initial ball of each
colour and replacement matrix $I_2$ is balanced. Its count law is
uniform. We prove that two separated closed tails of the colour
proportion have zero exponential decay at speed $n$. Consequently
no function $c(x-m)^2$ with $c>0$, for any centre $m$, can satisfy
the closed-set upper bound of a large-deviation principle. This
refutes the asserted implication from balance to positive
definiteness. The finite count recursion and the asymptotic
contradiction are both proved in Lean 4.
\end{abstract}

\section*{Model and scope}
At each step a ball is selected uniformly, returned, and one more
ball of its colour is added. Initially there is one red and one
blue ball. In the standard convention, the replacement matrix is
\[
 R=\begin{pmatrix}1&0\\0&1\end{pmatrix}.
\]
An urn is balanced when the two row sums are equal
\cite[Section 1]{bmw}. Here both sums equal $1$, and the total
number of balls after $n$ draws is $n+2$.

Let $K_n$ be the number of red draws and
$X_n=(K_n+1)/(n+2)$ the red proportion. We use the conventional
speed-$n$ large-deviation question for $X_n$. The conjecture imposes
no irreducibility or strictly positive off-diagonal hypothesis;
this standard balanced urn is therefore an admissible input.
The result refutes the positive-definiteness assertion under
balance. It does not assert that every balanced urn fails to have
a large-deviation principle, or that the zero quadratic form is
excluded by the first clause alone.

\begin{lemma}\label{lem:law}
For every $n\geq0$ and $0\leq k\leq n$,
$\mathbb P(K_n=k)=1/(n+1)$.
\end{lemma}
\begin{proof}
Write $p_n(k)=\mathbb P(K_n=k)$, with value zero outside
$0\leq k\leq n$. The initial law is $p_0(0)=1$.
For $1\leq k\leq n$, the two incoming transitions give
\[
 p_{n+1}(k)=p_n(k)\frac{n-k+1}{n+2}
             +p_n(k-1)\frac{k}{n+2}
           =\frac{1}{n+2}.
\]
At the two endpoints only one term survives:
$p_{n+1}(0)=p_n(0)(n+1)/(n+2)$ and
$p_{n+1}(n+1)=p_n(n)(n+1)/(n+2)$.
Both equal $1/(n+2)$. Induction proves the formula.
\end{proof}

\newpage
\section*{The large-deviation obstruction}
Set $F_-=[0,1/4]$ and $F_+=[3/4,1]$. For $n\geq2$,
the event $K_n=0$ implies $X_n\in F_-$, while $K_n=n$
implies $X_n\in F_+$. Lemma~\ref{lem:law} therefore gives
\[
 \frac{1}{n+1}\leq\mathbb P(X_n\in F_\pm)\leq1.
\]
Taking logarithms and dividing by $n>0$ yields
\[
 -\frac{\log(n+1)}{n}
 \leq\frac1n\log\mathbb P(X_n\in F_\pm)\leq0.
\]
Since $\log(n+1)/n\to0$, both middle expressions tend to zero.

\begin{theorem}
For every $c>0$ and $m\in\mathbb R$, the function
$I(x)=c(x-m)^2$ cannot be a speed-$n$ large-deviation rate
function for $X_n$, even when it is assigned $+\infty$ outside
$[0,1]$.
\end{theorem}
\begin{proof}
Such a rate would satisfy, for every closed set $F$,
\[
 \limsup_{n\to\infty}\frac1n\log\mathbb P(X_n\in F)
 \leq-\inf_{x\in F}I(x).
\]
If $m\leq1/2$, then $x-m\geq1/4$ on $F_+$, so
$\inf_{F_+}I\geq c/16>0$. The displayed upper bound says
$0\leq-c/16$, a contradiction. If $m>1/2$, then
$m-x\geq1/4$ on $F_-$, and the same contradiction follows
using $F_-$. Both sets are contained in $[0,1]$, so extending
the rate by infinity elsewhere has no effect on the argument.
\end{proof}

Thus balance does not imply a positive quadratic rate, even
after allowing an arbitrary centre. In particular the usual
candidate $c(x-1/2)^2$ is impossible. A positive definite
quadratic form for centred red/blue proportions restricts to
such a positive quadratic on their one-dimensional constraint
line, so changing to the two-coordinate description does not
remove this obstruction.

\section*{What is formally checked}
The Lean definition \texttt{mass} is the actual forward count
recursion, with the zero endpoint separated and natural
subtraction used only for the incoming blue term. Impossible
states have mass zero. The theorem \texttt{mass\_formula}
solves it for every $n,k$, and \texttt{mass\_total} proves
normalization. Event probabilities are finite sums of these
nonnegative masses at the actual proportion $(k+1)/(n+2)$;
\texttt{prob\_unit} proves that $[0,1]$ has probability one.

\newpage
\texttt{right\_log\_limit} and \texttt{left\_log\_limit}
prove the two real logarithmic limits, using the polynomial
lower bound and the logarithm's asymptotic growth. The definition
\texttt{ClosedUpperBound} writes the actual limsup/infimum
upper bound restricted to nonempty closed sets in $[0,1]$
with eventually positive probability. This is a necessary
part of the usual large-deviation upper bound. Restricting to
these sets avoids the zero-probability logarithm convention;
both sets used in the proof have strictly positive probability
for every $n\geq2$.

The theorem \texttt{no\_positive\_quadratic} proves that
this necessary condition fails for every $c>0$ and every real
$m$, using the actual infimum of the quadratic on the selected
closed interval. The final \texttt{counterexample} combines
that conclusion with the explicit replacement matrix's balance.
No existence theorem for a full LDP, or full classification of
all urn rate functions, is claimed or needed.

\section*{Reproduction and review}
The project pins Lean 4.19.0 and Mathlib revision\\
\texttt{c44e0c8ee63ca166450922a373c7409c5d26b00b}.
From the submitted \texttt{lean/} directory, run
\begin{verbatim}
lake exe cache get
lake build
lake env lean -DwarningAsError=true Main.lean
\end{verbatim}
The complete submission is freshly built. The axiom audit
contains only Lean's standard \texttt{propext},
\texttt{Classical.choice} and \texttt{Quot.sound}; there are no
proof placeholders, custom axioms or external computation
oracles.

The supplementary \texttt{verify.py} uses Python's standard
library and exact fractions. It reconstructs the Markov
transition law through $n=200$, and independently enumerates
all colour histories through $n=10$ with evolving ball counts.
It checks the uniform masses, normalization, support and the
two tail bounds. These finite checks are not used to justify
the asymptotic claim; the general formula and limits are
proved in Lean. Run \texttt{python verify.py}, and compile
this source with \texttt{tectonic main.tex}.

The exact bilingual conjecture is preserved in \texttt{SOURCE.md}.
Build, provenance and review records are included in
\texttt{verification/}. Review was an adversarial self-review
by the author; no independent reviewer or subagent was used.

\begin{thebibliography}{9}
\bibitem{source} The Last Math Competition,
Conjecture 00000002114. Exact source and upstream provenance
included with this submission.
\bibitem{bmw} C. Banderier, P. Marchal and M. Wallner,
\emph{Periodic P\'olya urns and an application to Young tableaux},
authors' preprint, February 7, 2018, Section 1.
\href{https://www.math.univ-paris13.fr/~marchal/aofa18.pdf}{Primary source}.
This reference supplies the standard urn and balance conventions;
the counterexample calculations above are self-contained.
\end{thebibliography}
\end{document}
